\documentclass[12pt,a4paper]{article}
\usepackage[utf8]{inputenc}
\usepackage[spanish]{babel}
\usepackage{geometry}
\geometry{margin=2.5cm}
\usepackage{enumitem}
\usepackage{titlesec}
\usepackage{array}
\usepackage{booktabs}
\usepackage{xcolor}
\usepackage{fancyhdr}
\usepackage{graphicx}
\usepackage{mdframed}
\usepackage{parskip}

\definecolor{azulUCT}{RGB}{0, 56, 116}
\definecolor{grisClaro}{RGB}{240, 240, 240}

\pagestyle{fancy}
\fancyhf{}
\fancyhead[L]{\small\textcolor{azulUCT}{\textbf{MAT1186 -- Introducción al Cálculo}}}
\fancyhead[R]{\small\textcolor{azulUCT}{Universidad Católica de Temuco}}
\fancyfoot[C]{\small\thepage}
\renewcommand{\headrulewidth}{0.4pt}
\renewcommand{\footrulewidth}{0.4pt}

\titleformat{\section}
  {\normalfont\Large\bfseries\color{azulUCT}}
  {\thesection}{1em}{}[\textcolor{azulUCT}{\titlerule[1pt]}]
\titleformat{\subsection}
  {\normalfont\large\bfseries\color{azulUCT!80}}
  {\thesubsection}{1em}{}
\titlespacing*{\section}{0pt}{1.5em}{0.8em}
\titlespacing*{\subsection}{0pt}{1em}{0.5em}

\newenvironment{articulo}[1]{%
  \begin{mdframed}[backgroundcolor=grisClaro,linecolor=azulUCT,linewidth=1.5pt,
    leftmargin=0pt,rightmargin=0pt,innerleftmargin=12pt,innerrightmargin=12pt,
    innertopmargin=8pt,innerbottommargin=8pt,roundcorner=4pt]
  \textbf{\textcolor{azulUCT}{Art. #1.}}\ %
}{\end{mdframed}\vspace{4pt}}

\begin{document}

\begin{titlepage}
  \centering
  \vspace*{1cm}
  \includegraphics[width=6cm]{logoUct.png}\par
  \vspace{0.6cm}
  {\large Ingeniería Civil en Informática -- Facultad de Ingeniería\par}
  \vspace{1.5cm}
  \textcolor{azulUCT}{\rule{\linewidth}{1.5pt}}
  \vspace{0.5cm}
  {\Huge\bfseries\textcolor{azulUCT}{Código de Ética}\par}
  \vspace{0.3cm}
  {\LARGE\bfseries y Reglamento Interno de Grupo\par}
  \vspace{0.5cm}
  \textcolor{azulUCT}{\rule{\linewidth}{1.5pt}}
  \vspace{1cm}
  {\large\textit{Evaluación Integrada de Desempeño N°1}\par}
  {\large MAT1186 -- Introducción al Cálculo\par}
  \vspace{2cm}
  \begin{tabular}{ll}
    \textbf{Líder de Grupo:} & Benjamín De La Fuente \\[6pt]
    \textbf{Suplente:}       & Sebastián Rodrigo Rivera Nahuelhuen \\[6pt]
    \textbf{Integrante:}     & Rodrigo Reyes \\
  \end{tabular}
  \vfill
  {\large Temuco, \today\par}
\end{titlepage}

\section*{Preámbulo}

El presente Código de Ética y Reglamento Interno regula el funcionamiento del grupo de trabajo
conformado para el desarrollo de la Evaluación Integrada de Desempeño N°1 del curso MAT1186.

Este documento surge del compromiso colectivo de sus integrantes con los valores de
\textbf{responsabilidad}, \textbf{honestidad académica}, \textbf{colaboración activa} y
\textbf{respeto mutuo}. Asumimos que el éxito del proyecto depende de la participación equitativa y
genuina de cada miembro, y que cualquier conducta contraria a esos principios perjudica tanto al
equipo como al proceso de aprendizaje individual.

Al suscribir este código, cada integrante declara haberlo leído, comprendido y acepta
voluntariamente su cumplimiento durante toda la duración del proyecto.

\newpage

\section{Identificación del Grupo}

\begin{center}
\renewcommand{\arraystretch}{1.6}
\begin{tabular}{>{\bfseries}p{4cm} p{8cm}}
  \toprule
  Curso      & MAT1186 -- Introducción al Cálculo \\
  Sección    & 1 \\
  Año / Sem. & 2026 -- Semestre 1 \\
  Proyecto   & Evaluación Integrada de Desempeño N°1 \\
  \midrule
  Líder      & Benjamín De La Fuente \\
  Suplente   & Sebastián Rodrigo Rivera Nahuelhuen \\
  Integrante & Rodrigo Reyes \\
  \bottomrule
\end{tabular}
\end{center}

\section{Definición de Roles y Responsabilidades}

\subsection{Líder de Grupo -- Benjamín De La Fuente}

\begin{itemize}[leftmargin=1.5em, itemsep=3pt]
  \item Organizar y convocar las reuniones de trabajo, fijando agenda y plazos internos.
  \item Ser el canal oficial de comunicación con el equipo docente.
  \item Supervisar el avance general del proyecto y velar por el cumplimiento de los plazos.
  \item Gestionar el repositorio de GitHub: crear la estructura inicial, administrar ramas y revisar que todos los integrantes registren sus commits.
  \item Mediar en conflictos internos antes de escalarlos al docente.
  \item Participar activamente en el desarrollo técnico del proyecto como un integrante más.
\end{itemize}

\subsection{Suplente -- Sebastián Rodrigo Rivera Nahuelhuen}

\begin{itemize}[leftmargin=1.5em, itemsep=3pt]
  \item Asumir las funciones de líder en ausencia justificada de este.
  \item Desarrollar los módulos técnicos asignados en coordinación con el equipo.
  \item Registrar avances mediante commits regulares y descriptivos en GitHub.
  \item Participar activamente en la preparación de la defensa oral.
  \item Informar oportunamente al líder si alguna tarea no puede completarse según lo acordado.
\end{itemize}

\subsection{Integrante -- Rodrigo Reyes}

\begin{itemize}[leftmargin=1.5em, itemsep=3pt]
  \item Desarrollar los módulos técnicos asignados según la distribución acordada por el equipo.
  \item Registrar avances mediante commits regulares y descriptivos en GitHub.
  \item Participar activamente en la preparación de la defensa oral.
  \item Documentar el código bajo los estándares acordados por el grupo.
  \item Informar oportunamente al líder si alguna tarea no puede completarse según lo acordado.
\end{itemize}

\section{Funcionamiento Interno}

\subsection{Reuniones y comunicación}

\begin{articulo}{1}
El grupo se reunirá con una frecuencia mínima de \textbf{dos veces por semana} durante el período de desarrollo del proyecto. Las reuniones podrán ser presenciales o en línea, según disponibilidad acordada.
\end{articulo}

\begin{articulo}{2}
El canal principal de comunicación será \textbf{WhatsApp}. Todos los integrantes deben responder mensajes en un plazo máximo de \textbf{24 horas}. Los acuerdos alcanzados en reuniones deberán quedar registrados por escrito en dicho canal.
\end{articulo}

\begin{articulo}{3}
La inasistencia a una reunión deberá ser comunicada con al menos \textbf{12 horas de anticipación}. La inasistencia reiterada sin justificación constituye una falta al presente código.
\end{articulo}

\subsection{Distribución de tareas}

\begin{articulo}{4}
Las tareas se distribuirán en la primera reunión formal del grupo, buscando una carga de trabajo equitativa. La distribución quedará registrada en un documento compartido y en el \textit{README} del repositorio.
\end{articulo}

\begin{articulo}{5}
Cada integrante es responsable de completar sus tareas asignadas en los plazos internos acordados. Si un miembro prevé dificultades, deberá informarlo con anticipación suficiente para que el equipo pueda reorganizarse.
\end{articulo}

\subsection{Repositorio y control de versiones}

\begin{articulo}{6}
Todos los integrantes deben realizar un mínimo de \textbf{tres (3) commits significativos} en el repositorio de GitHub. Los mensajes de commit deben ser claros y describir concretamente el cambio realizado; se prohíben mensajes genéricos como \texttt{update} o \texttt{fix}.
\end{articulo}

\begin{articulo}{7}
Queda prohibido realizar commits con trabajo ajeno como si fuera propio. Cada commit debe representar una contribución real del integrante que lo suscribe.
\end{articulo}

\section{Código de Ética y Conducta}

\subsection{Compromisos individuales}

Cada integrante del grupo declara y se compromete a:

\begin{enumerate}[leftmargin=1.5em, label=\textbf{\arabic*.}, itemsep=4pt]

  \item \textbf{Uso responsable de herramientas de apoyo:} el grupo podrá utilizar herramientas de inteligencia artificial como apoyo en el proceso de desarrollo. No obstante, todos los integrantes deben comprender, ser capaces de explicar y poder defender cualquier parte del trabajo resultante. Queda prohibida la entrega de código o contenido que ningún miembro del grupo comprende o puede justificar. La copia de repositorios externos sin atribución y cualquier otra forma de plagio constituyen faltas graves.

  \item \textbf{Participación activa y equitativa:} ningún integrante se limitará a observar el trabajo ajeno. Todos contribuirán al desarrollo técnico, matemático y a la preparación de la defensa oral.

  \item \textbf{Transparencia:} si un integrante no comprende algún módulo del proyecto, lo comunicará al grupo en lugar de ocultar la situación. El aprendizaje mutuo es parte del objetivo del trabajo.

  \item \textbf{Respeto:} las interacciones dentro del grupo se desarrollarán en un clima de respeto y tolerancia. Quedan prohibidas las descalificaciones personales, las críticas destructivas y toda forma de discriminación.

  \item \textbf{Puntualidad:} los integrantes se comprometen a cumplir los plazos internos acordados y a asistir puntualmente a las instancias de trabajo colectivo.

  \item \textbf{Responsabilidad colectiva:} todos los integrantes deben comprender el funcionamiento completo del programa. Nadie puede excusarse de responder en la defensa oral argumentando que ``esa parte la hizo otro''.

  \item \textbf{Integridad del repositorio:} el historial de commits refleja el trabajo real de cada integrante. Manipular el repositorio para simular participación que no existió constituye una falta grave.

\end{enumerate}

\subsection{Consecuencias por incumplimiento}

Las faltas al presente código se abordarán de manera escalonada:

\begin{center}
\renewcommand{\arraystretch}{1.5}
\begin{tabular}{>{\centering\bfseries\color{azulUCT}}m{2cm} >{\bfseries}m{3.5cm} m{7cm}}
  \toprule
  \textbf{Nivel} & \textbf{Tipo de falta} & \textbf{Medida} \\
  \midrule
  1 & Leve\newline(primera vez) &
    Llamado de atención verbal por parte del líder, con registro escrito en el canal del grupo. \\
  2 & Moderada\newline(reincidencia) &
    Reasignación de tareas y advertencia formal. Se documenta el incidente y se informa al docente de manera preventiva. \\
  3 & Grave\newline(reincidencia o falta ética mayor) &
    Comunicación formal al equipo docente, acompañada de evidencia, solicitando intervención externa. La calificación grupal podría verse afectada. \\
  \bottomrule
\end{tabular}
\end{center}

\section{Procedimiento ante Conflictos}

\begin{articulo}{8}
Ante cualquier desacuerdo, el grupo seguirá el siguiente protocolo:
\begin{enumerate}[leftmargin=1.2em, itemsep=2pt, label=\arabic{enumi}.]
  \item \textbf{Diálogo directo:} las partes involucradas intentarán resolver el conflicto de manera autónoma en el plazo de 48 horas.
  \item \textbf{Mediación del líder:} si no se alcanza un acuerdo, el líder actuará como mediador. Si el líder es parte del conflicto, asumirá esta función el suplente.
  \item \textbf{Escalamiento al docente:} si el conflicto persiste luego de la mediación interna, el grupo comunicará la situación al equipo docente, presentando el registro escrito de los intentos previos de resolución.
\end{enumerate}
\end{articulo}

\begin{articulo}{9}
Cualquier decisión que afecte la estructura o distribución del proyecto deberá ser consensuada por la mayoría del grupo (al menos 2 de 3 integrantes).
\end{articulo}

\section{Declaración Final}

Los abajo firmantes declaramos haber leído, comprendido y aceptado en su totalidad el presente Código de Ética y Reglamento Interno. Asumimos el compromiso de actuar conforme a sus principios durante el desarrollo de la Evaluación Integrada de Desempeño N°1 del curso MAT1186.

Reconocemos que el incumplimiento de este código puede afectar la calificación grupal y, principalmente, la integridad del aprendizaje que buscamos alcanzar como futuros profesionales.

\vspace{2.5cm}

\begin{center}
\renewcommand{\arraystretch}{2.5}
\begin{tabular}{ccc}
  \underline{\hspace{4.5cm}} & \underline{\hspace{4.5cm}} & \underline{\hspace{4.5cm}} \\[2pt]
  \textbf{Benjamín De La Fuente} & \textbf{Sebastián R. Rivera N.} & \textbf{Rodrigo Reyes} \\
  Líder de Grupo & Suplente & Integrante \\[20pt]
  Firma: \underline{\hspace{3cm}} & Firma: \underline{\hspace{3cm}} & Firma: \underline{\hspace{3cm}} \\[2pt]
  Fecha: \underline{\hspace{3cm}} & Fecha: \underline{\hspace{3cm}} & Fecha: \underline{\hspace{3cm}} \\
\end{tabular}
\end{center}

\end{document}